% ═══════════════════════════════════════════════════════════════════
% Section III: Proposed Scheduling Policy
% ═══════════════════════════════════════════════════════════════════

\section{Proposed Scheduling Policy}\label{Design}

% ── A. Continuous-Alpha Decision Rule ──────────────────────────────
\subsection{Continuous-Alpha Decision Rule}

Algorithm~\ref{alg:scheduler} summarizes the proposed policy.
At each slot, the monitor computes $\delta(t)$ from~\eqref{eq:delta} and maps it to a blending weight $\alpha(t)$ that interpolates continuously between full reuse and full LS reliance.

\begin{algorithm}[t]
\caption{Opportunistic CE Inference Scheduler}\label{alg:scheduler}
\begin{algorithmic}[1]
\Require LS estimates $\hat{\mathbf{H}}_{\text{LS}}(t)$, full-CE threshold $\tau_{\mathrm{full}}$, noise floor $\delta_{\min}$, safety counter $N_{\max}$
\State $n_{\text{skip}} \gets 0$
\For{each slot $t = 1, 2, \ldots, T$}
    \State Compute $\delta(t)$ via~\eqref{eq:delta}
    \If{$n_{\text{skip}} \geq N_{\max}$ \textbf{or} $\delta(t) > \tau_{\mathrm{full}}$} \Comment{Full CE}
        \State $a(t) \gets \text{Full}$;\; $n_{\text{skip}} \gets 0$
        \State $\hat{\mathbf{H}}(t) \gets \mathcal{E}\bigl(\hat{\mathbf{H}}_{\text{LS}}(t)\bigr)$
    \Else \Comment{Blend}
        \State $\alpha(t) \gets \mathrm{clamp}\!\Bigl(\frac{\delta(t) - \delta_{\min}}{\tau_{\mathrm{full}} - \delta_{\min}},\; 0,\; 1\Bigr)$ \Comment{via~\eqref{eq:alpha}}
        \State $\hat{\mathbf{H}}(t) \gets \bigl(1 - \alpha(t)\bigr)\,\hat{\mathbf{H}}(t{-}1) + \alpha(t)\,\hat{\mathbf{H}}_{\text{LS}}(t)$
        \State $n_{\text{skip}} \gets n_{\text{skip}} + 1$
    \EndIf
\EndFor
\end{algorithmic}
\end{algorithm}

The blending weight is given by
%
\begin{equation}\label{eq:alpha}
  \alpha(t) = \mathrm{clamp}\!\left(
    \frac{\delta(t) - \delta_{\min}}{\tau_{\mathrm{full}} - \delta_{\min}},\; 0,\; 1
  \right),
\end{equation}
%
where $\delta_{\min} = \sqrt{2/\gamma}$ is the LS noise floor at SNR $\gamma$.
When $\delta(t) \leq \delta_{\min}$, the observed variation is indistinguishable from estimation noise, so $\alpha(t) = 0$ and the previous estimate is reused exactly.
As $\delta(t)$ increases toward $\tau_{\mathrm{full}}$, $\alpha(t)$ grows linearly toward one, giving progressively more weight to the fresh LS observation.

The policy has two thresholds and one safety parameter.
The threshold $\tau_{\mathrm{full}}$ controls when the full CE method $\mathcal{E}$ is triggered.
The noise floor $\delta_{\min}$ (typically set to $\sqrt{2/\gamma}$) determines the boundary below which $\alpha(t) = 0$.
The safety counter $N_{\max}$ enforces a maximum number of consecutive blended slots, preventing unbounded staleness accumulation.

% ── B. Operating Modes ─────────────────────────────────────────────
\subsection{Operating Modes}

The scheduler operates in two modes.

\textbf{Blend} ($\delta \leq \tau_{\mathrm{full}}$).
The blending weight $\alpha(t)$ varies continuously from 0 (pure reuse) to 1 (full LS weight).
When $\delta(t)$ falls below the noise floor $\delta_{\min}$, $\alpha(t) = 0$ and the estimate is reused at zero additional cost beyond the LS monitor.
As $\delta(t)$ increases, more weight is given to the fresh LS observation, allowing the scheduler to track gradual channel drift without launching the full CE method.

\textbf{Full CE} ($\delta > \tau_{\mathrm{full}}$ or safety trigger).
A large channel change is detected, or the safety counter expires. The full CE method $\mathcal{E}$ is launched to produce a fresh high-quality estimate.

% ── C. CE-Algorithm Agnosticism ──────────────────────────────────
\subsection{CE-Algorithm Agnosticism}

The scheduler interfaces with the CE method $\mathcal{E}$ only through its output $\hat{\mathbf{H}}_{\mathcal{E}}(t)$.
It does not inspect the internal structure of $\mathcal{E}$, making it applicable to any CE algorithm.
We validate this property across three methods of increasing cost and accuracy.
%
\begin{itemize}
\item \textbf{LS} (baseline). $\mathcal{E} = \mathrm{identity}$; $c_{\mathcal{E}} \approx 0$.
\item \textbf{Genie-aided LMMSE}. Assumes perfect second-order statistics; $c_{\mathcal{E}} = \mathcal{O}(N_r^3 K)$.
\item \textbf{DL-CE (PACENet)}. Residual CNN with 57K parameters; $c_{\mathcal{E}} \approx 0.9$\,ms on A100 for ELAA.
\end{itemize}

% ── D. Evaluation Metrics ────────────────────────────────────────
\subsection{Skip-Aware Metrics}

We define four metrics that extend per-slot NMSE to the scheduling setting.

\textbf{Skip Rate (SR)} $= N_{\alpha=0} / T$, the fraction of slots where $\alpha(t) = 0$ (pure reuse).

\textbf{Blend Rate (BR)} $= N_{0 < \alpha < 1} / T$, the fraction of slots where $0 < \alpha(t) < 1$ (partial blending).

\textbf{Computation Ratio (CR)} $= \bar{c}_{\text{sched}} / c_{\mathcal{E}}$, the ratio of average per-slot cost under scheduling to the cost of always running $\mathcal{E}$.

\textbf{NMSE} is computed over all slots (including blended and reused ones) against the ground-truth channel, capturing the quality degradation from reuse and continuous blending.
